\documentclass[10pt,a4paper]{amsart}
\usepackage[margin=19mm]{geometry}
\usepackage{amsmath,amssymb,mathtools}
\usepackage[scr=rsfs,cal=euler]{mathalfa}
\usepackage{xfrac}
\usepackage[unicode,hidelinks,bookmarks=false]{hyperref}
\newcommand{\ie}{i.e.\ }
\newcommand{\cf}{cf.\ }
\newcommand{\resp}{respectively\ }
\newcommand{\Subsection}{\subsection}
\providecommand{\nobreakdash}{\nobreak\hbox{-}}
\setlength{\emergencystretch}{2em}
\multlinegap0pt
\usepackage[normalem]{ulem}
\usepackage{amssymb}
\usepackage[shortlabels]{enumitem}
\newtheorem{defn}{Definition}
\newtheorem{prop}{Proposition}
\title{A conformal reduction for the X-ADM mass}
\author{Stephen McCormick}
\date{Source: arXiv:2607.03629v1 (CC BY 4.0)}
\begin{document}
\maketitle

\noindent\emph{Source and licence.} A conformal reduction for the X-ADM mass. Version arXiv:2607.03629v1 (CC BY 4.0), distributed by arXiv under the Creative Commons Attribution 4.0 International licence, http://creativecommons.org/licenses/by/4.0/, as recorded in the arXiv licence metadata for this exact version and shown on the abstract page https://arxiv.org/abs/2607.03629v1. Full licence notice: \texttt{licenses/arxiv-2607.03629-CC-BY-4.0.txt}.\par\smallskip

\noindent\textit{Editorial scope.} Counted unit: the article's prop prop-Yamabe (source0.tex lines 254--260). The article's complete original proof of it is delivered with it (source0.tex lines 346--361, one \texttt{proof} environment of 16 lines, which the article places away from the statement). No mathematics is written, repaired or re-derived: the statement and the proof are the article's own lines, in the article's own notation, and no retained line is substituted or re-flowed.  Retained uncounted as the source-local context the proof needs: lines 222--224; lines 263--274.  Nothing was omitted from any retained span, so the package carries no omission record.  No retained line contains a citation, so the wrapper carries no bibliography and no bibliography entry.  No retained line contains a figure environment, an image-inclusion command or drawing code, so no image is delivered and the package declares no asset dependency.  Editorial formatting only, in addition: an AMS \texttt{amsart} preamble that restates the article's own declarations and macros used by the retained lines, namely the environments \texttt{defn}, \texttt{prop} on separate counters, together with the standard packages those lines need.  The retained environments print in this document as Definition 1, Proposition 1 in order of appearance, whereas the article prints the same items with the numbering of its own full document.  The complete native source of the article as distributed in the arXiv e-print for this exact version is bundled unchanged as \texttt{sources/204405/source0.tex}.
\begin{defn}\label{def-AF}
	We say $(M,g)$ is \textit{asymptotically flat} if $g-\mathring{g}\in W^{2,p}_\delta$ for some $p>\frac{n}{2}$ and $\delta\in(-(n-2),-\frac{n-2}{2}]$. \textit{We will assume $p$ and $\delta$ satisfy these restrictions throughout the entire document, for all weighted Sobolev and Lebesgue spaces used.}
\end{defn}

\begin{prop}\label{prop-coerce}
	Fix some one-form $A\in L^\infty_{-1-\varepsilon}$ and consider the operator
	\begin{equation*}
		Lu= du+Au.
	\end{equation*}
	
	For $u\in W^{1,2}_{\delta}$, we have
	\begin{equation}\label{eq-coerce}
		\|u\|_{2^*,\delta}\leq C\|Lu\|_{2,\delta-1},
	\end{equation}
	where $2^*=\frac{2n}{n-2}$ and $C$ depends only on $g$ and $A$.
\end{prop}

\begin{prop}\label{prop-Yamabe}
	Let $(M,g)$ be an asymptotically flat manifold in the sense of Definition \ref{def-AF} and let $X\in W^{2,p}_{\delta-1}$ be a vector field on $M$. Suppose
	\begin{equation}
		R(g)+ 2\mathrm{div}(X)\geq\frac{n-2}{n-1}|X|^2,
	\end{equation}
	then $g$ is positive Yamabe type.
\end{prop}

\begin{proof}[Proof of Proposition \ref{prop-Yamabe}]
	By direct computation we have for any $f\in C^\infty_c$,
	\begin{align*}
		\int_M &\frac{4(n-1)}{n-2}|\nabla f|^2+R(g)f^2\,dV_g\geq \\
		 &\geq \int_M \frac{4(n-1)}{n-2}|\nabla f|^2-\left( 2\mathrm{div}(X)-\frac{n-2}{n-1}|X|^2 \right)f^2\,dV\\
	&=\int_M \frac{4(n-1)}{n-2}|\nabla f|^2+4fX\cdot\nabla f+\frac{n-2}{n-1}|X|^2f^2\,dV\\
	&=\int_M \frac{4(n-1)}{n-2}\left|\nabla f +\frac{n-2}{2(n-1)}Xf   \right|^2\,dV.
	\end{align*}

	Now note that $W^{2,p}_{\delta-1}\subset L^\infty_{-1-\varepsilon}$ so Proposition \ref{prop-coerce} applies to the operator $L=d+\frac{n-2}{2(n-1)}X^\flat$. In particular, we apply Proposition \ref{prop-coerce} with $\delta=-(n-2)/2$ and then since the weighted space $L^{2^*}_{-(n-2)/2}$ is equal to the unweighted space $L^{2^*}$, we have
	\begin{equation*}
		\|f\|^2_{2^*}\leq C \int_M \frac{4(n-1)}{n-2}\left|\nabla f +\frac{n-2}{2(n-1)}Xf   \right|^2\,dV.
	\end{equation*}
	From this, we immediately see $\mathcal Y([g])>0$.
	
\end{proof}

\end{document}
